\documentclass[11pt,a4paper]{article}
\usepackage[margin=21mm,headheight=15pt]{geometry}
\usepackage[T1]{fontenc}
\usepackage{lmodern,amsmath,amssymb,booktabs,tabularx,enumitem,xcolor,fancyhdr}
\usepackage[hidelinks]{hyperref}
\definecolor{accent}{RGB}{0,114,178}
\setlength{\parindent}{0pt}
\setlength{\parskip}{6pt}
\setlist{nosep,leftmargin=*}
\pagestyle{fancy}
\fancyhf{}
\fancyhead[L]{\small Econometrics 25117 | Instructor guide}
\fancyhead[R]{\small 2026--2027}
\fancyfoot[L]{\small Private teaching notes}
\fancyfoot[R]{\thepage}
\newcommand{\heading}[1]{\section*{\color{accent}#1}}
\newcommand{\cue}[2]{\par\noindent\textbf{#1}\quad #2\par}
\setlength{\emergencystretch}{1em}
\newcommand{\E}{\mathbb{E}}
\newcommand{\Var}{\operatorname{Var}}
\newcommand{\Cov}{\operatorname{Cov}}
\begin{document}
\begin{center}{\Large\bfseries Midterm 2 instructor guide}\end{center}
\textbf{Slide route:} Syllabus date: November 18, 2026. Proposed coverage: Lectures 5--7, with earlier prerequisites.
\heading{Session 10 | Midterm 2}
\textbf{Purpose:} Second assessment session; no new content. The syllabus gives November 18. Calendar placement and exam duration require confirmation against the actual timetable.
\heading{Proposed blueprint}
\begin{tabularx}{\textwidth}{@{}l r X@{}}\toprule
Area & Share & Evidence of learning\\\midrule
Multiple regression & 35\% & OVB sign, partial interpretation, reference categories.\\
Joint inference & 30\% & Joint null, restrictions, covariance-aware test.\\
Functional forms & 35\% & Log units, quadratic effect, interaction contrast.\\\bottomrule
\end{tabularx}
The blueprint is a planning aid, not an approved exam. Announce final coverage in advance. If the course requires cumulative coverage for this midterm, adjust the blueprint explicitly.
\cue{100-minute slot assumption}{Ten minutes for instructions, 80 for the paper, ten for collection; replace with official timing if different. Confirm approved resources, room, and accessibility arrangements in advance.}
\cue{What students should know}{A regression coefficient needs units and a conditioning statement. A joint test needs a joint null. An interaction effect needs a specified value of the interacting variable.}
\heading{Instructor calibration examples}
\cue{OVB}{An omitted positive determinant of $Y$ negatively correlated with $X$ creates negative bias in a simple-regression slope under the usual linear setup. Award credit for both links in the argument.}
\cue{Restrictions}{Testing $\beta_1=\beta_2$ is one restriction, not two. The variance of their estimated difference includes their covariance.}
\cue{Interactions}{In $Y=5+2X+3D+4XD+u$, the group contrast at $X=2$ is $3+8=11$. The slope difference is 4. These answer different questions.}
\cue{Logs}{A log-outcome coefficient $.1$ means an approximate 10\% change, with exact exponentiated-index change $100(e^{.1}-1)\simeq10.52\%$.}
\heading{After the paper}
Record which errors concern algebra and which concern interpretation. Prepare a short return-to-teaching activity using a strong regression result with a doubtful causal story. This provides the transition to threats to identification.
\cue{Next session}{Lecture8v2.tex. Ask: even with correct standard errors and a flexible functional form, what can still go wrong?}

\end{document}
